\section{Empirical implications}\label{sec:empirical}

The model gives predictions for publicly visible sale processes in which one bidder is prepared and a challenger has
not yet committed to diligence. Each prediction names the incumbent, prior, trading-pattern and tie-rule conditions
it needs, and concerns one target in one public state; Appendix~\ref{app:empirical} states each in full and proves
it. The predictions concern price impacts, the informed investor's presence, the winner and announcement returns.
Under full two-sided trade, class \(H\) buys and class \(L\) sells, each with probability one, and the informed share
of nonzero orders is \(u/(2+u)\) at every prior. In the example at the prior \(1/2\) and the example's slot, full
two-sided trade is the unique outcome at every strength between \(r_S\), the strength at which a full sale earns the
trading cost, and \(h\); this range contains \(r_1\), the entry ceiling \(r_C\) and \(r_2\).

\paragraph{Entry and incumbent strength.} When the price reaches the challenger before its decision, a stronger
prepared bidder can raise the probability that the challenger prepares; in the example, that probability rises by
\(\reg{entry_jump_B}\) from \(r_0\) to \(r_1\). When the price does not reach the challenger, a stronger prepared
bidder can only lower it. Against a strong enough prepared bidder, entry can return to its floor while the price stays
informative. Proposition~\ref{prop:ccc} gives the first and third comparisons, and the fixed-information control of
Section~\ref{sec:benchmark} gives the second (Proposition~\ref{prop:control}).

\paragraph{The price gap and its entry component.} Under full two-sided trade at a prior below the preparation
threshold, the purchase impact \(P(1)-P(0)\), weighted by the frequency of a purchase, exceeds the sale impact
\(P(0)-P(-1)\), weighted by the frequency of a sale, exactly when a purchase recruits the challenger
(Proposition~\ref{prop:spread}). The weighted difference equals the expected price before the order opportunity
minus the price after no order, so the test needs prices and order frequencies and no auction object; at the prior
\(1/2\) the two frequencies are equal, and the test compares the two impacts directly. Under the same conditions the
price gap \(P(1)-P(-1)\) has a class component, \(\rho\DT\,[M(\pi)-m(\pi)]\), and an entry component, which is
positive exactly when a purchase recruits; the whole entry component sits in the purchase impact. As the incumbent
strengthens, entry adds more to the price through competition and less by rescuing a sale that the incumbent alone
would not complete. With a uniform incumbent and the prior and the slot fixed, the entry component rises with
strength while full two-sided trade and recruitment persist, when \(3p^2<\min\{\underline\theta,\underline r\}^2\),
where \(\underline\theta\) is the lowest challenger value and \(\underline r\) the lowest strength compared; the
example meets this condition.

In the example at \(r_1\), the price gap is \(\reg{ba_r1_B}\), and its entry component is \(\reg{ba_share_r1_B}\) of
it. With a hidden price, the price gap would be \(\reg{ba_class_r1_B}\). At \(r_2\) the price gap is
\(\reg{ba_r2_B}\), the class component alone, and at the entry ceiling \(r_C\) it falls by \(\reg{ba_jump_rC_B}\).

\paragraph{The informed share.} Under full two-sided trade the mean order is \(\lambda(2\pi-1)\), so order counts
identify the informed investor's presence when the prior is known and differs from one half
(Proposition~\ref{prop:informedshare}). The sale impact identifies presence at any known prior, one half included:
when entry after a sale and after no order sits at its floor and the payoff scale \(\rho\DT\) is known, the sale
impact weighted by the frequency of a sale, divided by \(\rho\DT\), equals \(\pi(1-\pi)\lambda\), and together with
the mean order \(\lambda(2\pi-1)\) it identifies both the presence probability and the prior. At the prior
\(1/2\) the counts depend on the slot only through the volume of nonzero orders, so prices separate slots that counts
cannot. In the example at \(r_1\), the shares consistent with the counts and the investor's incentives form the
interval from \(\approx\reg{theta_min_counts_B}\) to \(\reg{theta_max_full_eq_r1_B}\), and a slot with share
\(\reg{alt_theta_B}\) has the example's volume and entry, but its price gap at \(r_1\) is \(\reg{alt_ba_r1_B}\),
against \(\reg{ba_r1_B}\).

\paragraph{Winners.} Under full two-sided trade, a prepared challenger wins more often after a purchase than after a
sale exactly when the incumbent's value falls, with positive probability, where a high-class value beats it with a
higher probability than a low-class value; the gap is \(M(\pi)-m(\pi)\) times the difference between the classes' chances of
beating the incumbent (Proposition~\ref{prop:winner}). Orders predict the winner even where they do not predict
entry: where entry stays at its floor after every order, as above the entry ceiling, the order still covaries
positively with a challenger win under the same condition, and not with entry. With the prior and the slot fixed and
full two-sided trade at each strength compared, when class \(H\) outbids every incumbent the gap does not narrow as
the incumbent strengthens, and it widens wherever the stronger incumbent takes wins from the low class. Among
prepared challengers, the share of the high class at each price is the price-implied posterior, which rises with the
price under full two-sided trade. In the example class \(H\) outbids every incumbent, and the gap is
\(\reg{wgap_r1_B}\) at \(r_1\) and \(\reg{wgap_r2_B}\) at \(r_2\); the covariance of the order with entry is
\(\reg{cov_entry_r1_B}\) at \(r_1\) and zero at \(r_2\).

\paragraph{Announcements.} For a positive target-payoff spread and \(g_H>g_L\), if a public event after the price reveals only whether the
challenger prepared, the mean price change over preparation reports is positive exactly when preparation was still uncertain after some order, at
every prior and in every equilibrium (Proposition~\ref{prop:announcement}); in the example it is
\(\reg{evprep_mean_r1_B}\) at \(r_1\). The mean change over both reports, with their model probabilities, is zero. Under
full two-sided trade at a prior below the preparation threshold, an opportunity that passes with no order is no news
about the class and, when a purchase recruits, bad news about entry: in the example at \(r_1\), the price after no
order minus the expected price before the opportunity is \(\reg{pdrop_neutral_r1_B}\), and the size of this drop is
the weighted difference of the price-impact test.

\paragraph{Two clusters at intermediate strengths.} In the example, at one public strength in the range
\((r_F,r_N)\) of Section~\ref{sec:coexistence}, processes can form two clusters. In one cluster the market is quiet:
the mean order is zero and the three prices coincide. In the other, the high class buys and the low class abstains;
the mean order is then \(\lambda/2\) (\(\reg{imb_buy_B}\) in the example), every challenger prepares after a purchase,
and a sale executes above the price after no order, by \(\reg{U_gap_r160_B}\) at \(\reg{str_r160}\).

\paragraph{Measurement.} The counterpart of preparation is the start of substantive diligence or a proposal that
requires costly preparation; public offers are a narrower outcome. Incumbent strength is measured from information
that is public before that decision, and the price must precede it.

The order is the first prespecified public order opportunity after the sale becomes public, and an opportunity with
no order is recorded as a zero. A scheduled call auction is one example, such as the opening auction on the first
trading day after the approach is disclosed; its time is fixed in advance and does not depend on any trader's
information, and the auction takes place whether or not an order arrives. The sign of its net order imbalance is the
order, and an imbalance below a size fixed in advance records the zero order. A sample of first executed trades omits
the zero order, so it cannot identify the volume, the effect of an inactive opportunity, or the informed share.

The announcement formulas need events that reveal only preparation. Target returns before a bid also
partly reflect anticipation of that bid \citep{Schwert1996}, so each test dates the price before the challenger's
decision. The trading counterpart is outside trade with a view on the challenger's fit, such as orders from sector
specialists.
